\section{Why an angular bound alone has no converse in finite order}
\label{sec:analytic-angle-only-obstruction}

\begin{theorem}[NN-A026: fixed nonreal angles with arbitrarily high finite positivity]
For every integer $n\ge1$ put
\[
 \vartheta_n=\frac{\pi}{3n},\qquad
 g_n(z)=\prod_{j=0}^{n}
                (1+e^{i\vartheta_n(n-2j)}z),\qquad
 Q(z)=\cosh\sqrt z,\qquad F_n(z)=g_n(z)Q(z).
\]
Then $F_n$ is real entire of order $1/2$, $F_n(0)=1$, and all
its Taylor coefficients are strictly positive. Every allowed
Toeplitz minor of order at most $2n$ is strictly positive.
Nevertheless $F_n$ has the two fixed nonreal zeros
$-e^{i\pi/3}$ and $-e^{-i\pi/3}$, and its maximal zero-angle
deviation from the negative real ray is exactly $\pi/3$.
Its least negative-minor order is finite and exceeds $2n$.

Rescaling to $\widehat F_n(z)=F_n(z/F_n'(0))$ also imposes
$\widehat F_n'(0)=1$ and preserves the same angular and finite
positivity conclusions. Separately, put
$\Psi_n(z)=F_n(z)/F_n(1)$. Then
$U\mapsto\det\Psi_n(U)$ is a normalized central positive-type
function on every $U(d)$ with $d\le2n$, with strictly positive
coefficients for every polynomial irreducible, yet fails to be
positive type at some larger finite rank.

Thus the maximal zero angle, the entire growth order $1/2$,
and the displayed coefficient normalizations alone give no
upper bound on necessary detecting order or unitary rank.
The positive lower bounds in NN-A022 use an angular sufficient
condition; a matching upper bound for the actual heat family
requires additional information about that family.
\end{theorem}

\begin{proof}
Let $m=n+1$, $k=2n$, and more generally consider the centered
geometric alphabet
\[
 x_j=e^{i\vartheta(m+1-2j)},\quad 1\le j\le m,
                \qquad 0<\vartheta\le\frac{\pi}{m+k-1}.
\]
For a partition $\lambda$ with at most $m$ rows, the Weyl
alternant evaluated at this geometric progression gives
\begin{equation}\label{eq:analytic-geometric-schur-sines}
 s_\lambda(x_1,\ldots,x_m)
     =\prod_{1\le i<j\le m}
        \frac{\sin(\vartheta(\lambda_i-\lambda_j+j-i))}
             {\sin(\vartheta(j-i))}.
\end{equation}
For completeness, write the alternant columns as geometric
powers with exponents $\lambda_j+m-j$. Both determinants
become Vandermonde products. Replacing
$e^{iu}-e^{-iu}$ by $2i\sin u$ gives the quotient above;
the overall phases cancel because the alphabet is centered.
The denominator is nonzero, since $0<\vartheta(j-i)<\pi$.
The formula also follows directly from the Weyl character
formula, Etingof et al., Theorem~5.22.1, printed p.~132.

If $\lambda_1\le k$, every numerator argument lies in $[0,\pi]$:
\[
 1\le\lambda_i-\lambda_j+j-i\le k+m-1.
\]
Hence all ordinary Schur values of width at most $k$ are
nonnegative. Values with more than $m$ rows are zero by the
alternating definition. Every skew Schur function of width
at most $k$ is a nonnegative integer linear combination of
ordinary Schur functions of that width: the restriction
multiplicities used and proved in NN-A021 are nonnegative,
and a nonzero coefficient entails containment of the relevant
partition in the outer partition. The exact minor/width
correspondence in NN-A020 now proves that the coefficients
of $\prod_j(1+x_jz)$ are $PF_k$.

Taking $m=n+1,k=2n$ and
$\vartheta=\pi/(m+k-1)=\pi/(3n)$ gives the stated $g_n$.
Its roots are $-x_j^{-1}$; the endpoint phases are exactly
$\pm\pi/3$, with all other phases between them. Its
coefficients are nonnegative by the already proved order-one
minors. Conjugate pairing of the alphabet makes $g_n$ real.

The entire function $Q(z)=\cosh\sqrt z$ has Taylor coefficients
$1/(2r)!$, and the genus-zero product
\[
 Q(z)=\prod_{j\ge1}\left(1+
                       \frac{4z}{\pi^2(2j-1)^2}\right).
\]
Thus $Q$ is $PF_\infty$ and has strictly positive allowed
Toeplitz minors, by the explicit positive-path proof in
NN-A008. Its coefficients are strictly positive. Multiplying
by $g_n$ preserves $PF_{2n}$ by Cauchy--Binet. It actually
preserves strict positivity of every allowed minor: in
$T_QT_{g_n}$, select the intermediate index set equal to the
column index set of the target minor. The selected
$T_{g_n}$ minor is upper triangular with diagonal one,
while the selected $T_Q$ minor is the strictly positive
allowed minor. Every other Cauchy--Binet term is nonnegative.
The sum is finite for each fixed index set, by upper
triangularity. The coefficients of the product are positive
because the constant coefficient of $g_n$ is one.

A polynomial factor changes neither the order $1/2$ of $Q$
nor its negative-real infinite zero set. The fixed nonreal
zeros of $g_n$ remain, so the maximal angle is exactly
$\pi/3$. Edrei's entire-function converse, with the exact
original locator in NN-A007, rules out $PF_\infty$;
therefore some finite minor is negative. None has order
at most $2n$. Positive scaling of the variable multiplies
any Toeplitz minor by a positive number and preserves angles,
which proves the derivative-normalized assertion.

Finally $F_n(1)>0$, so $\Psi_n(0)>0$ and $\Psi_n(1)=1$.
The exact $PF_d$ versus positive-$U(d)$ criterion in NN-A021
proves the asserted characters, including strict positivity
of all polynomial-irreducible coefficients. The same
criterion and Edrei give failure at some finite larger rank.
This character normalization and the derivative normalization
are separate operations; no preservation under a negative
translation of the variable is assumed.
\end{proof}

\paragraph{Classical locator and the new use.}
The finite sharp family is the Schoenberg example recorded in
Adm--Garloff--Tyaglov, \emph{Total nonnegativity of finite
Hurwitz matrices and root location of polynomials},
Theorem~5.7 and equation~(5.9), printed p.~161, immediately before the sector
sufficiency Theorem~5.8, equation~(5.10). The displayed Weyl
calculation supplies a complete proof of precisely the
finite-order assertion needed here. The order-$1/2$ infinite
extension, strict-minor observation, normalizations and
finite-unitary-rank interpretation are given above. This is
an obstruction to an unrestricted converse, not a counterexample
to a heat-specific upper estimate.

\paragraph{An explicit correction to the source's product inference.}
Immediately after equations~(5.11)--(5.12), printed p.~162,
the same published source asserts that a product of quadratics
whose angles are exterior to the quadratic $PF_k$ interval
cannot belong to $PF_k$. That inference is false, even when
every factor has the same outside angle. Set
\[
 q(z)=1+\frac9{10}z+z^2,\qquad
 \theta=\arccos\frac9{20}>\frac\pi3.
\]
Then $q=(1+e^{i\theta}z)(1+e^{-i\theta}z)$ is not $PF_2$,
since its central order-two minor is $81/100-1<0$.
On the other hand
\[
 q(z)^2=1+\frac95z+\frac{281}{100}z^2+\frac95z^3+z^4
\]
is $PF_2$. Its three nontrivial adjacent log-concavity
differences are $43/100$, $46561/10000$, and $43/100$,
all positive. For a positive finite consecutive coefficient
sequence, these differences make the successive ratios
nonincreasing. Therefore every order-two Toeplitz minor
\[
       a_{h+d}a_{h+e}-a_ha_{h+d+e}\quad(d,e>0)
\]
is nonnegative by multiplying the corresponding ratio
inequalities; when an index is outside the support,
upper triangularity and the consecutive support give the
same conclusion directly. This proves the required $PF_2$
assertion without a numerical test. The exact published
source is
\url{https://www-home.htwg-konstanz.de/~garloff/JMAA-3_Adm_Garloff_Tyaglov.pdf},
DOI~10.1016/j.jmaa.2018.06.065. The valid sector sufficiency
and valid single-quadratic criterion used in NN-A020 are
unaffected; they were also proved there independently.

\paragraph{Geometric and representation meaning.}
The values in \eqref{eq:analytic-geometric-schur-sines} are
literal Weyl characters evaluated at a centered diagonal
unitary matrix. In the stated rectangle of partition widths,
they are nonnegative real numbers despite its nonreal
eigenvalues. Their skew restrictions give the positive
Toeplitz tests, and the independent positive $Q$ specialization
makes all allowed tests strict. The representation-ring
exterior-power and exact-sequence formalism has the exact
locator SGA~6, Expos\'e~V, Example~2.2(b) and Lemma~2.2.1.
The functions on each $U(d)$ are actual mixtures of its
polynomial irreducible characters; they fail to extend
positively through all ranks for the fixed analytic function.
Their real coefficients are not integral representation
multiplicities. The example separates finite character-cone
positivity from a constraint depending only on the support
angle of the analytic zero divisor.
